% STATUS: WRITTEN, 15 November 2026, after the Evaluation section, so that it
% describes the result that was measured rather than the one that was hoped
% for.

\section{Conclusion and Future Work}

This paper asked whether a reinforcement learning agent can act as an
autonomous database administrator, re-fragmenting and re-allocating shards
as a workload's geographic center of gravity moves. What it establishes is a
measured problem statement, a cluster that confirms the premise, a learned
policy that makes a different and cheaper trade than the obvious heuristic,
and evidence that the agent's place among the policies carries over from
the simulator to the cluster.

On a real six-node CockroachDB cluster, gathering leaseholders in the
issuing region rather than spreading them changes p95 latency by
1337.8\,ms. A reactive heuristic that follows the traffic is 16.9\,ms above
an exact per-window oracle under slow drift and 75.4\,ms above it under fast
drift,
so the case for a learned policy is conditional on volatility, and this work
measures the condition.

At volatility 4 the agent is 13.7\% faster than static hashing on held-out
seeds and faster than a Lion-style co-location planner on all four learner
seeds, and on average it moves about a fifth of the data the heuristic
moves. It is 17\,ms
slower than the heuristic, and at the migration weight it was trained with
the heuristic still has the lower objective cost. Rescoring the same
behavior shows the verdict turns on the price of movement: in simulation the
agent is the cheapest policy for any migration weight between 0.49 and
2.18. On the cluster, over four workload seeds, the heuristic is cheapest and
the agent second on every seed, 12.6\% faster than static hashing on
average, and the simulator is optimistic by a steady 14\% for static
hashing and the heuristic. The Lion-style planner lost more in transfer,
partly to lease moves the database refused onto replicas that had not yet
caught up, a rule the simulator does not model.

% ---------------------------------------------------------------------
\subsection{Future work}

In the order they should be attempted. Add learner seeds to the cluster
evaluation, so it has an interval, and train an agent with splitting
disabled, to measure what fragmentation adds. Train at the migration weight a
deployment cares about, which the rescoring in Table~\ref{tab:weight} can
suggest but not answer. Teach the simulator the database's refusal rule, or
retry a refused move in the next window. Address the plateau with the other
Rainbow components~\cite{hessel2018rainbow}, and choose the checkpoint on
more evaluation episodes and on the objective rather than latency alone.
Finally, fine-tune on the cluster: the harness already executes actions,
reads placement back and computes the simulator's reward, and at about 15\,s
a step the budget, not the code, is the constraint.
